\documentclass[hyperref={pdfpagelabels=false},table]{beamer}

\input{lec_common}
\date{}
\newcommand{\fontsmall}{\fontsize{5pt}{6pt}\selectfont}
\begin{document}

\part{Unit 10: Pandas}

%\frame{\tableofcontents}

\begin{frame}[fragile]{A Guiding Example}
This unit needs an example data source to demonstrate
Pandas:  solar cells in Södra Sandby\\[0.5cm]
\begin{minipage}[h]{0.4\textwidth}
\lstset{basicstyle=\tiny}
\pyth!solarWatts.dat!
\begin{lstlisting}
Date;Watt
           :
2019-10-20 08:22:00 ; 44.0
2019-10-20 08:23:00 ; 61.0
2019-10-20 08:24:00 ; 42.0
           :
\end{lstlisting}
\end{minipage}
\begin{minipage}[h]{0.4\textwidth}
\lstset{basicstyle=\tiny}
\pyth!price.dat!
\begin{lstlisting}
Date;SEK
           :
2019-10-20 01:00 ; 0.32
2019-10-20 02:00 ; 0.28
2019-10-20 03:00 ; 0.29
           :
\end{lstlisting}
\end{minipage}
\begin{minipage}[h]{0.4\textwidth}
\lstset{basicstyle=\tiny}
\pyth!rates.dat!
\begin{lstlisting}
Date;Euro_SEK
       :
2019-10-21 ; 10.7311
2019-10-22 ; 10.7303
2019-10-23 ; 10.7385
       :

\end{lstlisting}
\end{minipage}
\end{frame}
\begin{frame}[fragile]{Typical questions?}
\begin{itemize}
\item When was the highest power production?
\item Which day delivered the most energy?
\item Which was the economical outcome in Euro per month?
\item Was there a day where the solar cells were out of work?
\end{itemize}
\end{frame}
\begin{frame}[fragile]
\vfill
The tool is Pandas dataframe
\vfill
\end{frame}
\begin{frame}[fragile]{Pandas Dataframe versus Numpy Array}
\vfill
Let's start by just looking at an example
\vfill
\begin{minipage}{0.45\textwidth}
\begin{python}
A=array(
     [[1.,2.,3.],
      [4.,5.,6.]])
print(A)
\end{python}
results in
\begin{verbatim}
[[1. 2. 3.]
 [4. 5. 6.]]
\end{verbatim}
\end{minipage}
\begin{minipage}{0.45\textwidth}
\begin{python}
import pandas as pd
AF = pd.DataFrame(A)
print(AF)
\end{python}
results in
\begin{verbatim}
     0    1    2
0  1.0  2.0  3.0
1  4.0  5.0  6.0
\end{verbatim}
\end{minipage}
\vfill
We see that a Pandas dataframe has extra labels of the rows and columns - called
\emph{index} and \emph{columns}. These are metadata of a dataframe.

Here they coincide with NumPy's indexing.

\end{frame}
\begin{frame}[fragile]{Dataframe: Index and Column Metadata}

Index and Column Metadata gives a Pandas dataframe the possibility to label
the data as known from classical tabel design

\begin{python}
AF.columns = ['C1','C2','C3']
AF.index    = ['R1', 'R2']
\end{python}
\vfill
This gives
\begin{verbatim}
     C1   C2   C3
R1  1.0  2.0  3.0
R2  4.0  5.0  6.0
\end{verbatim}
\vfill
... and allows  indexing to construct subframes:
\vfill
\begin{minipage}{0.45\textwidth}
\begin{python}
AF.loc[['R1'],['C2']]
\end{python}
\end{minipage} \hfill
returns \hfill
\begin{minipage}{0.25\textwidth}
\begin{verbatim}
     C2
R1  2.0
\end{verbatim}
\end{minipage}
\end{frame}
\begin{frame}[fragile]{Dataframe: Indexing}
Indexing by row and column labels \pyth!loc!-method:
\vfill
\begin{minipage}{0.45\textwidth}
\begin{python}
AF.loc['R1', 'C2']     # returns 2.0
\end{python}
\end{minipage} \vfill
while \vfill
\begin{minipage}{0.45\textwidth}
\begin{python}
AF.loc[['R1'],['C2']]
\end{python}
\end{minipage} \hfill
returns \hfill
\begin{minipage}{0.25\textwidth}
\begin{verbatim}
     C2
R1  2.0
\end{verbatim}
\end{minipage}
\end{frame}
\begin{frame}[fragile]{Indexing by row and column indexes \pyth!iloc!-method}
\vfill
\begin{minipage}{0.45\textwidth}
\begin{python}
AF.iloc[0,1]     # returns 2.0
\end{python}
\end{minipage} \vfill
while \vfill
\begin{minipage}{0.45\textwidth}
\begin{python}
AF.iloc[[0],[1]]
\end{python}
\end{minipage} \hfill
returns \hfill
\begin{minipage}{0.25\textwidth}
\begin{verbatim}
     C2
R1   2.0
\end{verbatim}
\end{minipage}
\end{frame}
\begin{frame}[fragile]{Pandas Dataseries}
The label of a column can be used as an attribute
\begin{python}
AF.C1
\end{python}
which refers to the entire column
\begin{verbatim}
R1    1.0
R2    4.0
Name: C1, dtype: float64
\end{verbatim}
\vfill
This is a new Pandas object -- a Series object
\begin{python}
type(AF.C1) == pd.Series  # True
\end{python}
Note, a Pandas series has no column label. It is just a single column corresponding
to a single type of measured data.
\end{frame}
\begin{frame}[fragile]{Dataframes: Subframes, Slices and Elements}
Note, if \pyth!loc! (or \pyth!iloc!) is called with list arguments or slices  the result is
a dataframe
\begin{python}
AF.loc['R1':,'C2':]  # or equivalently
AF.loc[['R1','R2'], ['C2','C3']]
\end{python}
\vfill
while calling with a pair of single labels just give an element of the dataframe
\begin{python}
AF.loc['R1','C2']  # returns 2.0
\end{python}
\vfill
{\tiny Compare with indexing of arrays}
\end{frame}
\begin{frame}[fragile]{Datetime objects as indexes}
It is very common that the index contains information about the date and the time of a given measurement
as a string.

In Sweden we use a format \texttt{yy-mm-dd hh:mm:ss} to report the time of a measurement as a string.

\begin{python}
date_index=['2020-01-19 11:45', '2020-01-20 02:14']
AF.index=date_index
\end{python}
\vfill
results in the dataframe
\vfill
\begin{verbatim}
                  C1   C2   C3
2020-01-19 11:45  1.0  2.0  3.0
2020-01-20 02:14  4.0  5.0  6.0
\end{verbatim}
\end{frame}
\begin{frame}[fragile]{Datetime objects as indexes \hfill (Cont.)}
We often want to make simple arithmetic with date and time information for this end
the list of strings has to be converted in a list of so-called \pyth!datetime! objects first
\vfill
\begin{python}
date_index=pd.to_datetime(
           ['2020-01-19 11:45', '2020-01-20 02:14']
           )
AF.index=date_index
\end{python}
\vfill
This way we can easlily find out how much time passed between the measurements
\pyth!AF.index[1]-AF.index[0]!.
The result is a \pyth!timedelta! object displayed here as
\begin{verbatim}
Timedelta('0 days 14:29:00')
\end{verbatim}
\end{frame}
\begin{frame}[fragile]{Creating a dataframe from imported data}
A dataframe from solar data (File: \pyth!solarWatts.dat!)
\end{frame}
\begin{frame}[fragile]{Dataframe from a datafile}
\begin{python}
solarWatts = pd.read_csv("solarWatts.dat",
                         sep=';',
                         index_col='Date',
                         parse_dates=[0],
                         infer_datetime_format=True)
\end{python}
What do these arguments tell?
The data is semicolon seperated, which column defines the index, which column contains dates, should dates converted from string to timestamp and how.
\vfill
A typical entry: \hfill \pyth!solarWatts.iloc[100]!\\
The last date of the measurements: \hfill \pyth!solarWatts.iloc[-1]!
\end{frame}
\begin{frame}[fragile]{Merging Dataframes}
Here we merge production data with electricity prices:
\begin{python}
# Creating dataframes from csv-files
solarWatts = pd.read_csv("solarWatts.dat", ... )
prices     = pd.read_csv("price.dat", ...)
rates      = pd.read_csv("rates.dat", ...)
\end{python}
and
\begin{python}
# Merging dataframes
solar_all=pd.merge(solarWatts, prices, how='outer', sort=True, on='Date')
solar_all=pd.merge(solar_all, rates, how='outer', sort=True, on='Date')
\end{python}
\pyth!'outer'! tells that the union of the index entries is taken as a new index.
\end{frame}
\begin{frame}[fragile]{Missing data}
In the example:\\ power is updated every minute, price is updated every hour, rate is updated every day\\
\vfill
This gives enties like:
\begin{lstlisting}
                        Watt      SEK  Euro_SEK
Date
2019-10-21 13:59:00  2650.0      NaN       NaN
2019-10-21 14:00:00  1901.0  0.48725       NaN
2019-10-21 14:01:00  1436.0      NaN       NaN
2019-10-21 14:02:00  1040.0      NaN       NaN
2019-10-21 14:03:00   714.0      NaN       NaN
2019-10-21 14:04:00   498.0      NaN       NaN
\end{lstlisting}
What can be done? \\ Use frame methods like: \pyth!dropna!, \pyth!fillna!, \pyth!interpolate!
\end{frame}
\begin{frame}[fragile]{Missing data \hfill Cont.}
We use here  \pyth!fillna!:
\begin{python}
solar_all=solar_all.fillna(method='pad',axis=0)
\end{python}
Note, this method is no \emph{inplace} operation, i.e. a new object is created.\\
\vfill
\begin{lstlisting}
                       Watt      SEK  Euro_SEK
Date
2019-10-21 13:59:00  2650.0  0.47572   10.7311
2019-10-21 14:00:00  1901.0  0.48725   10.7311
2019-10-21 14:01:00  1436.0  0.48725   10.7311
2019-10-21 14:02:00  1040.0  0.48725   10.7311
2019-10-21 14:03:00   714.0  0.48725   10.7311
2019-10-21 14:04:00   498.0  0.48725   10.7311
\end{lstlisting}
\end{frame}
\begin{frame}[fragile]{Plotting dataframes}
We demonstrate this by plotting the price variation along one day.
\begin{python}
solar_all.loc['2020-05-16'].plot(y='SEK')
\end{python}
\centerline{
\includegraphics[width=.8\textheight]{SEK_plot1.png}}
\end{frame}
\end{document}
